\documentclass[11pt,a4paper]{article}

\usepackage{lmodern}
\usepackage[T1]{fontenc}
\usepackage[utf8]{inputenc}
\usepackage[margin=2.2cm]{geometry}
\usepackage{amsmath,amssymb}
\usepackage{graphicx}
\usepackage{booktabs}
\usepackage{listings}
\usepackage{xcolor}
\usepackage[font=small,labelfont=bf,skip=4pt]{caption}
\usepackage{subcaption}
\usepackage[hidelinks]{hyperref}
\usepackage{microtype}
\emergencystretch=3em   % long \texttt tokens must not overflow the margin

\graphicspath{{figures/}{../results/figures/}}

\definecolor{kw}{RGB}{0,90,160}
\definecolor{cm}{RGB}{90,120,90}
\definecolor{st}{RGB}{150,60,20}
\lstset{
  basicstyle=\ttfamily\scriptsize,
  keywordstyle=\color{kw}\bfseries,
  commentstyle=\color{cm}\itshape,
  stringstyle=\color{st},
  numbers=left, numberstyle=\tiny\color{gray}, numbersep=6pt,
  frame=single, rulecolor=\color{gray!40},
  breaklines=true, showstringspaces=false, captionpos=b,
  language=Python, tabsize=4, xleftmargin=14pt,
  % listings lit le fichier octet par octet : sans ces deux lignes, un accent dans
  % un listing casse la compilation avec « Invalid UTF-8 byte sequence ».
  extendedchars=true,
  literate={é}{{\'e}}1 {è}{{\`e}}1 {ê}{{\^e}}1 {ë}{{\"e}}1 {É}{{\'E}}1
           {à}{{\`a}}1 {â}{{\^a}}1 {ù}{{\`u}}1 {û}{{\^u}}1 {ü}{{\"u}}1
           {ô}{{\^o}}1 {î}{{\^\i}}1 {ï}{{\"\i}}1 {ç}{{\c c}}1 {œ}{{\oe}}1,
}

\newcommand{\AlsBestHR}{0.1186}
\newcommand{\AlsTimeToBest}{17.5}
\newcommand{\ConvItersAmz}{12}
\newcommand{\ConvItersMl}{17}
\newcommand{\CorrWorst}{7\cdot 10^{-14}}
\newcommand{\CrossKAmz}{34}
\newcommand{\CrossKMl}{213}
\newcommand{\EalsOverAls}{35}
\newcommand{\EalsOverAlsBest}{33}
\newcommand{\EalsOverAlsMl}{9}
\newcommand{\EalsTimeToAls}{3.8}
\newcommand{\FitConstAlsAmz}{0.23}
\newcommand{\FitConstAlsMl}{0.16}
\newcommand{\FitConstEalsAmz}{0.47}
\newcommand{\FitConstEalsMl}{0.46}
\newcommand{\FitCrossAlsAmz}{78}
\newcommand{\FitCrossAlsMl}{843}
\newcommand{\FitCrossEalsAmz}{315}
\newcommand{\FitCrossEalsMl}{1045}
\newcommand{\FitRtwoAlsAmz}{0.9999}
\newcommand{\FitRtwoAlsMl}{1.0000}
\newcommand{\FitRtwoEalsAmz}{0.9996}
\newcommand{\FitRtwoEalsMl}{0.9995}
\newcommand{\FlipKAmz}{128}
\newcommand{\FlipKMl}{128}
\newcommand{\GrowthAlsAmz}{10.3}
\newcommand{\GrowthAlsMl}{7.6}
\newcommand{\GrowthEalsAmz}{2.2}
\newcommand{\GrowthEalsMl}{1.8}
\newcommand{\HRa}{0.1171}
\newcommand{\HRaMl}{0.3373}
\newcommand{\HRe}{0.1578}
\newcommand{\HReMl}{0.3674}
\newcommand{\HRu}{0.1434}
\newcommand{\HRuMl}{0.3387}
\newcommand{\KAtRatioAmz}{128}
\newcommand{\KAtRatioMinAmz}{16}
\newcommand{\KAtRatioMinMl}{16}
\newcommand{\KAtRatioMl}{128}
\newcommand{\PopFloor}{3.1}
\newcommand{\PopFloorMl}{2.4}
\newcommand{\PopGain}{10.0}
\newcommand{\PopGainMl}{8.5}
\newcommand{\RatioAtMaxAmz}{2.7}
\newcommand{\RatioAtMaxMl}{1.7}
\newcommand{\RatioAtMinAmz}{2.2}
\newcommand{\RatioAtMinMl}{2.8}
\newcommand{\SerialBig}{0.378}
\newcommand{\SerialSmall}{0.385}
\newcommand{\SlopeAls}{1.68}
\newcommand{\SlopeAlsAmz}{1.30}
\newcommand{\SlopeAlsMl}{1.13}
\newcommand{\SlopeEals}{0.86}
\newcommand{\SlopeEalsAmz}{0.47}
\newcommand{\SlopeEalsMl}{0.40}
\newcommand{\SlopeHiAlsAmz}{1.68}
\newcommand{\SlopeHiAlsMl}{1.47}
\newcommand{\SlopeHiEalsAmz}{0.86}
\newcommand{\SlopeHiEalsMl}{0.66}
\newcommand{\SpeedupBig}{2.07}
\newcommand{\SpeedupSmall}{1.84}
\newcommand{\TeightBig}{6.72}
\newcommand{\TeightSmall}{0.64}
\newcommand{\ToneBig}{13.91}
\newcommand{\ToneSmall}{1.17}


% \input adds a \relax after the file, which opens a spurious cell when the
% fragment sits between \midrule and \bottomrule. The TeX primitive does not.
\makeatletter\newcommand{\inputrows}[1]{\@@input #1 }\makeatother

% A figure whose experiment has not been run must not break the build.
\newcommand{\figplot}[2]{%
  \IfFileExists{../results/figures/#1.pdf}%
    {\includegraphics[width=#2]{#1}}%
    {\fbox{\parbox{#2}{\centering\vspace{2.5em}\emph{La figure
      \texttt{\detokenize{#1}} n'a pas été générée : l'expérience
      correspondante n'a pas été exécutée.}\vspace{2.5em}}}}}

\title{\vspace{-1.2cm}\textbf{Factorisation matricielle rapide pour la recommandation en ligne\\
à partir de retours implicites : une implémentation Spark et\\ une étude expérimentale de sa scalabilité}}
\author{Thomas Sinapi \and Anaël Ernadote\\[4pt]
\normalsize M2 IASD - Université Paris-Dauphine PSL\\
\normalsize Projet \emph{Machine Learning for Big Data} - 2026}
\date{}

\begin{document}
\maketitle
\vspace{-0.8cm}

\begin{abstract}
\noindent
He et al.~\cite{he2016eals} proposent \textbf{eALS}, un algorithme de recommandation
qui apprend à partir de simples traces d'usage (clics, achats) plutôt que de notes.
Sa particularité est de traiter les articles que l'utilisateur n'a pas vus non pas
tous de la même façon, mais en fonction de leur popularité. Les auteurs affirment que
leur algorithme se parallélise très facilement, sans jamais le vérifier : ils ne
fournissent aucune implémentation distribuée.

C'est ce que fait ce rapport. Nous écrivons eALS en PySpark, ainsi qu'une version
séquentielle de référence en NumPy, et nous montrons que les deux produisent
exactement le même modèle (écart maximal $\CorrWorst$) : la parallélisation ne coûte
aucune précision. Nous mesurons ensuite le comportement du code sur un portable à
8 cœurs et quatre jeux de données de $10^5$ à $10^7$ interactions, dont Amazon Movies,
l'un des deux jeux utilisés par l'article.

Deux résultats principaux. Côté précision, eALS obtient $+\EalsOverAls$\,\% de HR@100
par rapport à l'ALS de Spark MLlib, et son coût croît beaucoup plus lentement avec le
nombre de facteurs $K$ ($\GrowthEalsAmz\times$ contre $\GrowthAlsAmz\times$ entre
$K=32$ et $K=128$). Côté parallélisme en revanche, l'accélération plafonne à
$\SpeedupBig\times$ sur 8 cœurs : la limite n'est pas l'algorithme, mais la machine.
\end{abstract}

\vspace{0.2cm}
\tableofcontents
\vspace{0.3cm}

%==========================================================================
\section{Introduction et énoncé du problème}
%==========================================================================

\paragraph{Le problème.} Un système de recommandation apprend rarement à partir de
notes. Il apprend à partir de traces : ce que l'utilisateur a cliqué, lu, acheté. Ces
traces sont abondantes, mais elles n'ont qu'un seul côté. Elles disent ce qu'un
utilisateur a fait, jamais ce qu'il a rejeté. Un article non consulté peut aussi bien
signifier « ça ne m'intéresse pas » que « je ne l'ai jamais vu passer ».

C'est une vraie difficulté, pas un détail. Si l'on entraîne un modèle uniquement sur
les interactions observées, il apprend à répondre « oui » partout, et ce modèle absurde
obtient un score parfait : il n'a jamais vu un seul exemple négatif. La solution
classique, due à Hu et al.~\cite{hu2008implicit}, consiste à considérer que toute paire
(utilisateur, article) non observée est un négatif, mais un négatif \emph{faible},
auquel on accorde un poids $w_0$ petit et identique partout.

\paragraph{Ce que change eALS.} L'article~\cite{he2016eals} conteste ce poids unique.
Ne pas avoir vu un film à succès est un indice assez fort de désintérêt ; ne pas avoir
vu un film confidentiel n'indique presque rien. Les auteurs remplacent donc le poids
unique $w_0$ par un poids $c_i$ propre à chaque article, croissant avec sa popularité.

Le problème est que ce changement, apparemment anodin, casse l'algorithme existant.
L'ALS classique calcule un vecteur entier de préférences à la fois, ce qui l'oblige à
inverser une matrice pour chaque utilisateur et chaque article. Cette opération n'est
abordable que grâce à un précalcul partagé, et ce précalcul n'est possible que si le
poids est identique partout. Autrement dit : dès que l'on veut des poids différenciés,
la recette ne fonctionne plus.

La réponse de He et al.\ est de changer de granularité. Plutôt que de résoudre un
vecteur entier, ils mettent à jour \emph{une seule coordonnée à la fois}. Chaque mise à
jour devient alors une simple division, sans aucune inversion de matrice. C'est ce que
signifie le « e » d'eALS : \emph{element-wise}, élément par élément. Le gain est d'un
facteur $K$ (le nombre de dimensions du modèle) sur le coût d'une itération, et les
poids peuvent enfin varier d'un article à l'autre.

\paragraph{Le point de départ de ce projet.} Ces mises à jour ne lisent que des données
partagées en lecture seule. Les auteurs en concluent qu'eALS est « trivialement
parallélisable » : on peut confier des groupes d'utilisateurs différents à des machines
différentes, et le résultat sera \emph{exactement} celui du calcul séquentiel, pas une
approximation.

L'article s'arrête là. Aucune implémentation distribuée n'est fournie, et cette
affirmation n'est jamais vérifiée. C'est précisément ce que demande le sujet du projet,
et ce que fait ce rapport : écrire l'implémentation Spark manquante, vérifier que le
parallélisme est bien exact, et mesurer ce qu'il rapporte réellement.

\paragraph{Plan.} La Section~\ref{sec:solution} décrit la solution retenue et justifie
les choix de conception. La Section~\ref{sec:algo} explique l'algorithme et comment
nous le distribuons. La Section~\ref{sec:code} commente les fragments de code qui
comptent. La Section~\ref{sec:exp} présente les mesures. La
Section~\ref{sec:discussion} discute forces et faiblesses. Les annexes contiennent le
code complet et un notebook de démonstration.

%==========================================================================
\section{La solution retenue}\label{sec:solution}
%==========================================================================

\subsection{Architecture générale}

Le système est un petit paquet Python plat (Table~\ref{tab:modules}). Il n'y a
délibérément ni framework, ni hiérarchie de classes, ni couche de configuration : ce
qui est intéressant, c'est le noyau numérique et la façon dont Spark l'ordonnance, et
les deux doivent tenir sur un écran.

\begin{table}[h]
\centering\small
\caption{Modules. Tout est en Python ; NumPy fait l'algèbre locale, Spark fait la
distribution.}\label{tab:modules}
\begin{tabular}{@{}lll@{}}
\toprule
Module & Rôle & Lignes \\
\midrule
\texttt{spark\_utils.py}      & l'unique endroit où la \texttt{SparkSession} est configurée & 42 \\
\texttt{data.py}              & chargement, filtrage 10-core, ré-indexation dense, découpage LOO & 133 \\
\texttt{eals\_local.py}       & référence NumPy \emph{et} les deux noyaux de mise à jour partagés & 147 \\
\texttt{eals\_rdd.py}         & implémentation Spark RDD (fondée sur les broadcasts) & 173 \\
\texttt{baselines.py}         & ALS implicite de MLlib~\cite{meng2016mllib}, MostPopular & 37 \\
\texttt{evaluate.py}          & classement sur le catalogue entier, HR@$k$, NDCG@$k$ & 69 \\
\texttt{check\_correctness.py}& les deux implémentations doivent concorder & 99 \\
\texttt{experiments.py}       & chaque mesure $\rightarrow$ \texttt{results/*.csv} & 217 \\
\texttt{plots.py}             & chaque figure $\rightarrow$ \texttt{results/figures} & 158 \\
\bottomrule
\end{tabular}
\end{table}

Une propriété que nous jugeons importante : \textbf{le noyau numérique n'existe qu'une
seule fois}. \texttt{eals\_local.update\_users} et \texttt{update\_items} sont
importées telles quelles par la version RDD. Seul l'\emph{ordonnancement} diffère entre
les deux. Tout écart entre elles localise donc un bug de plomberie, jamais un bug de
formule - et la Section~\ref{sec:correctness} montre qu'il n'y en a aucun.

\subsection{RDD ou DataFrame ?}

L'énoncé autorise l'un ou l'autre. Nous avons retenu les
\textbf{RDD~\cite{zaharia2012rdd} avec broadcast}, et le choix se justifie par un
argument de volume de communication plutôt que par une préférence d'API. Les deux
parallélisations envisageables ne sont pas deux écritures d'un même programme :

\begin{itemize}\itemsep2pt
\item \textbf{RDD + broadcast} (retenu). Les matrices de facteurs vivent sur le driver ;
à chaque étape, celle qui est en lecture seule est diffusée (broadcast), et chaque tâche
met à jour les lignes qu'elle possède. Par étape, le réseau transporte $O(NK)$ (le
broadcast) plus $O(MK)$ (le résultat collecté).
\item \textbf{DataFrame + jointure/\texttt{cogroup}}. Les
facteurs restent distribués sous forme de DataFrames ; une tâche reçoit les facteurs
dont elle a besoin via une jointure avec shuffle. Rien n'est jamais assemblé sur le
driver, donc le modèle peut dépasser sa mémoire - mais le shuffle transporte
$O(|\mathcal{R}|K)$ par étape au lieu de $O((M+N)K)$. Sur ml-1m ce rapport vaut
$|\mathcal{R}|/N \approx 300$.
\end{itemize}

Sur la machine visée - un seul portable, un modèle de quelques dizaines de Mo - la
contrainte de mémoire du driver qui justifierait la seconde option ne se présente
jamais, tandis que le facteur $\approx 300$ sur le volume échangé, lui, se paie à chaque
itération. Nous implémentons donc uniquement la version RDD. La limite de cette
décision est discutée en Section~\ref{sec:discussion} : c'est elle qui borne la taille
de modèle que ce code peut traiter.

\subsection{Choix du jeu de données}\label{sec:dataset}

Nous utilisons deux familles de jeux de données, pour deux usages distincts.

\textbf{Amazon Movies} est l'un des deux jeux de l'article, et c'est là que nous
mesurons la précision. Il est très creux (99,87\,\% de cases vides), ce qui correspond
au régime d'un vrai catalogue. Nous partons du fichier de notes de la catégorie
Films \& TV du graphe produits Amazon~\cite{mcauley2015amazon}, et non du dump brut cité
par l'article, qui n'est plus commode à obtenir ; c'est la même source et le même
domaine, sur un instantané un peu différent. Nos chiffres de précision ne sont donc pas
directement comparables à ceux de l'article, et nous ne le prétendons pas.

\textbf{MovieLens 100k / 1M / 10M}~\cite{harper2015movielens} sert aux mesures de temps. Une étude de scalabilité
a besoin du même jeu de données à plusieurs tailles, ce qu'Amazon ne fournit pas :
MovieLens offre une échelle propre en $\times100$ avec une sémantique identique à chaque
palier.

Les quatre jeux subissent le même prétraitement : les notes sont ignorées (une
interaction vaut $1$), on ne garde que les utilisateurs et les articles ayant au moins
10 interactions (filtrage \emph{10-core}), les identifiants sont renumérotés de façon
contiguë, et l'on met de côté la dernière interaction de chaque utilisateur pour le
test (protocole \emph{leave-one-out}).

\begin{table}[h]
\centering\small
\caption{Jeux de données après filtrage 10-core et ré-indexation dense.}\label{tab:datasets}
\begin{tabular}{@{}lrrrrrr@{}}
\toprule
Jeu de données & Utilisateurs $M$ & Items $N$ & $|\mathcal{R}|$ & Sparsité & Train & Test LOO \\
\midrule
\inputrows{tables/tab_datasets.tex}
\bottomrule
\end{tabular}
\end{table}

\subsection{Matériel et logiciels}\label{sec:setup}

Tout tourne sur un unique \textbf{ordinateur portable Apple série M, 8 cœurs
physiques, 16\,Go}, avec \textbf{Spark 4.1.0} en mode \texttt{local[$p$]},
\textbf{Python 3.11}, \textbf{NumPy 2.3} et \textbf{OpenJDK 17}. Il n'y a pas de
cluster : $p$, le nombre de cœurs, est le réglage que font varier les expériences de
parallélisme.

Un point de configuration mérite d'être signalé car il conditionne la validité des
mesures. Chaque tâche Spark appelle NumPy, et NumPy lance de lui-même ses propres
threads de calcul. Sur 8 cœurs, cela donne 8 tâches $\times$ 8 threads $=$ 64 threads en
concurrence, et la mesure à « 1 cœur » se retrouve silencieusement parallélisée : toute
accélération calculée ensuite est dénuée de sens. Nous bridons donc NumPy à un seul
thread (\texttt{scripts/env.sh}), de sorte que Spark soit la seule source de
parallélisme. Signalons aussi que Spark 4 exige Java 17 ou 21 ; sur un JDK plus récent,
aucun job ne démarre.

Tous les temps rapportés sont des médianes. Chaque configuration est répétée 3 fois,
chaque exécution effectue 4 itérations dont la première est écartée (échauffement), soit
9 mesures par point ; les barres d'erreur des figures donnent le minimum et le maximum.
La graine aléatoire est fixée partout, si bien que deux exécutions identiques
produisent le même modèle au bit près.

%==========================================================================
\section{L'algorithme}\label{sec:algo}
%==========================================================================

\subsection{Modèle et objectif}

\paragraph{L'idée du modèle.} On dispose d'un tableau géant à $M$ lignes
(les utilisateurs) et $N$ colonnes (les articles), rempli de $1$ là où une interaction a
eu lieu et de $0$ partout ailleurs. Ce tableau est immense et presque vide : sur Amazon
Movies, plus de $99{,}8$\,\% des cases sont à zéro.

La factorisation matricielle consiste à décrire chaque utilisateur et chaque article par
une courte liste de $K$ nombres, appelés \emph{facteurs latents} - typiquement $K=32$.
On peut se les représenter comme des goûts non nommés : « films d'action », « cinéma
d'auteur », etc., découverts automatiquement plutôt que définis à la main. La note
prédite pour un couple (utilisateur, article) est alors simplement le produit scalaire
de leurs deux listes : si les goûts coïncident, le score est élevé.

Apprendre le modèle revient donc à chercher les listes de nombres qui reproduisent au
mieux le tableau observé. Formellement, avec $P$ la matrice des facteurs utilisateurs,
$Q$ celle des articles et $\hat r_{ui}=\mathbf{p}_u^{\top}\mathbf{q}_i$ la prédiction,
on minimise l'écart quadratique sur \emph{toutes} les cases, chacune pondérée par un
poids $w_{ui}$ (Éq.~2 de l'article) :
\begin{equation}\label{eq:obj}
J \;=\; \sum_{u=1}^{M}\sum_{i=1}^{N} w_{ui}\,(r_{ui}-\hat r_{ui})^2
\;+\;\lambda\Big(\textstyle\sum_u \|\mathbf p_u\|^2+\sum_i\|\mathbf q_i\|^2\Big),
\end{equation}
avec $r_{ui}=1,\ w_{ui}=1$ sur les paires observées et $r_{ui}=0,\ w_{ui}=c_i$ ailleurs.
Le second terme, piloté par $\lambda$, est une régularisation : il pénalise les
facteurs de grande amplitude et empêche le modèle de surapprendre. En séparant les
cases observées des cases vides, on obtient la forme réellement optimisée (Éq.~7) :
\begin{equation}\label{eq:obj7}
J=\sum_{(u,i)\in\mathcal{R}} w_{ui}(r_{ui}-\hat r_{ui})^2
+\sum_{u=1}^{M}\sum_{i\notin\mathcal{R}_u} c_i\,\hat r_{ui}^{\,2}
+\lambda\Big(\textstyle\sum_u\|\mathbf p_u\|^2+\sum_i\|\mathbf q_i\|^2\Big).
\end{equation}

\paragraph{Le poids des articles non vus (Éq.~8).} C'est ici que se situe l'apport de
l'article. Le poids accordé à une case vide de la colonne $i$ vaut
\begin{equation}\label{eq:ci}
c_i \;=\; c_0\,\frac{f_i^{\alpha}}{\sum_{j=1}^{N} f_j^{\alpha}},
\qquad f_i=\frac{|\mathcal{R}_i|}{\sum_{j}|\mathcal{R}_j|},
\end{equation}
où $f_i$ est la fréquence de l'article $i$. La formule a une propriété commode : quelle
que soit la valeur de $\alpha$, la somme des $c_i$ vaut toujours $c_0$. Les deux
paramètres ont donc des rôles bien séparés. $c_0$ fixe l'importance \emph{totale}
accordée aux cases vides, et $\alpha$ décide seulement de la façon de la répartir entre
les articles.

Deux cas limites éclairent le mécanisme. Avec $\alpha=0$, tous les articles reçoivent le
même poids et l'on retrouve exactement la méthode de Hu et al. Avec $\alpha>0$, les
articles populaires en reçoivent davantage : ne pas avoir vu un film à succès pèse plus
lourd dans l'apprentissage que ne pas avoir vu un film obscur. Sur Amazon Movies, le
rapport entre le plus fort et le plus faible poids passe de $1$ à $\alpha=0$ à environ
$16$ à $\alpha=0{,}4$, la valeur retenue par l'article. Cette redistribution fait tout
l'intérêt de la méthode, et aussi, comme nous le verrons en
Section~\ref{sec:discussion}, tout son risque.

\subsection{Dériver la mise à jour élément par élément}

Cette sous-section refait le calcul qui mène à la règle de mise à jour. Elle mérite
d'être suivie, car c'est elle qui fait apparaître l'astuce sur laquelle repose toute
l'efficacité de la méthode - et, plus loin, toute la parallélisation.

Le principe est simple : on fixe tous les paramètres sauf un, et on cherche la valeur
optimale de celui-là. Comme la fonction à minimiser est un polynôme du second degré en
ce paramètre, il suffit d'annuler la dérivée pour obtenir directement l'optimum. Pas de
pas d'apprentissage, pas de tâtonnement.

Notons $\hat r^{\,f}_{ui}=\hat r_{ui}-p_{uf}q_{if}$ la prédiction \emph{privée de la
contribution du facteur $f$}. En dérivant~\eqref{eq:obj} par rapport à la seule
coordonnée $p_{uf}$ :
\begin{equation}
\frac{\partial J}{\partial p_{uf}}
= -2\sum_{i=1}^{N}\big(r_{ui}-\hat r^{\,f}_{ui}\big)w_{ui}q_{if}
+ 2\,p_{uf}\sum_{i=1}^{N} w_{ui}q_{if}^2 + 2\lambda p_{uf}.
\end{equation}
En l'annulant on obtient la solution élément par élément générique (Éq.~5)
\begin{equation}\label{eq:eq5}
p_{uf}=\frac{\sum_{i=1}^{N}\big(r_{ui}-\hat r^{\,f}_{ui}\big)w_{ui}q_{if}}
             {\sum_{i=1}^{N} w_{ui}q_{if}^2+\lambda},
\end{equation}
Cette formule est exacte, mais inutilisable telle quelle : les sommes portent sur les
$N$ articles du catalogue, pour \emph{chaque} utilisateur et \emph{chaque} facteur. On
parcourrait donc tout le tableau vide à chaque itération, soit $O(MNK)$ opérations. Sur
Amazon Movies, cela représente $23$ milliards de termes par itération, alors que seules
$10^6$ cases contiennent une information.

Toute l'astuce consiste à ne jamais parcourir ces cases vides. Séparons la somme en deux
parties, les cases observées et les autres, en utilisant $r_{ui}=0$ et $w_{ui}=c_i$ sur
les secondes :
\begin{equation}\label{eq:split}
p_{uf}=\frac{\sum_{i\in\mathcal{R}_u}(r_{ui}-\hat r^{\,f}_{ui})w_{ui}q_{if}
             \;-\;\sum_{i\notin\mathcal{R}_u}\hat r^{\,f}_{ui}\,c_i q_{if}}
            {\sum_{i\in\mathcal{R}_u} w_{ui}q_{if}^2
             +\sum_{i\notin\mathcal{R}_u} c_i q_{if}^2+\lambda}.
\end{equation}
Les deux sommes sur les cases vides ($i\notin\mathcal{R}_u$) restent le goulot
d'étranglement. L'idée décisive est de les réécrire par complément : \emph{tous les
articles, moins ceux que l'utilisateur a vus}. Le premier terme ne dépend alors plus de
l'utilisateur, et le second ne porte que sur les quelques articles réellement observés.
Pour le numérateur, en développant $\hat r^{\,f}_{ui}=\sum_{k\neq f}p_{uk}q_{ik}$ :
\begin{align}
\sum_{i\notin\mathcal{R}_u}\hat r^{\,f}_{ui}c_iq_{if}
&=\sum_{i=1}^{N} c_iq_{if}\sum_{k\neq f}p_{uk}q_{ik}
 -\sum_{i\in\mathcal{R}_u}\hat r^{\,f}_{ui}c_iq_{if} \nonumber\\
&=\sum_{k\neq f}p_{uk}\underbrace{\sum_{i=1}^{N} c_i q_{if}q_{ik}}_{\textstyle s^q_{kf}}
 -\sum_{i\in\mathcal{R}_u}\hat r^{\,f}_{ui}c_iq_{if},
\label{eq:num}
\end{align}
et pour le dénominateur, $\sum_{i\notin\mathcal{R}_u}c_iq_{if}^2
= s^q_{ff}-\sum_{i\in\mathcal{R}_u}c_iq_{if}^2$. L'observation cruciale est que la
matrice $K\times K$
\begin{equation}\label{eq:sq}
S^q \;=\; \sum_{i=1}^{N} c_i\,\mathbf q_i\mathbf q_i^{\top}
\end{equation}
ne dépend pas de l'utilisateur $u$. C'est le point clé de toute la méthode : cette
petite matrice $K\times K$ (soit $32\times32$ en pratique) résume à elle seule la
contribution des millions de cases vides. On la calcule \emph{une seule fois} par
itération, puis on la réutilise pour les $M$ utilisateurs. Le coût des données
manquantes passe ainsi de « proportionnel au catalogue » à « constant ».

En reportant~\eqref{eq:num} dans~\eqref{eq:split}, on obtient la règle effectivement
implémentée :
\begin{equation}\label{eq:pu}
\boxed{\;p_{uf}\leftarrow
\frac{\displaystyle\sum_{i\in\mathcal{R}_u}\!\big[w_{ui}r_{ui}-(w_{ui}-c_i)\hat r^{\,f}_{ui}\big]q_{if}
-\sum_{k\neq f}p_{uk}\,s^q_{kf}}
{\displaystyle\sum_{i\in\mathcal{R}_u}(w_{ui}-c_i)q_{if}^2+s^q_{ff}+\lambda}\;}
\end{equation}
Aucune somme ne porte plus sur les $N$ articles : elles portent sur les interactions
réelles de l'utilisateur, ou sur les $K$ facteurs. C'est ce qui rend la méthode
utilisable.

\paragraph{Côté article.} Le calcul est symétrique, avec une différence utile : le poids
$c_i$ ne dépend que de l'article, il peut donc sortir de la somme. Avec
$S^p=P^{\top}P$ le second cache (non pondéré, cette fois) :
\begin{equation}\label{eq:qi}
\boxed{\;q_{if}\leftarrow
\frac{\displaystyle\sum_{u\in\mathcal{R}_i}\!\big[w_{ui}r_{ui}-(w_{ui}-c_i)\hat r^{\,f}_{ui}\big]p_{uf}
-c_i\sum_{k\neq f}q_{ik}\,s^p_{kf}}
{\displaystyle\sum_{u\in\mathcal{R}_i}(w_{ui}-c_i)p_{uf}^2+c_i\,s^p_{ff}+\lambda}\;}
\end{equation}
C'est l'Éq.~(13) de l'article.

Cette réécriture est une optimisation, et une optimisation peut être fausse. Avant
d'écrire la moindre ligne de Spark, nous avons donc comparé les deux règles encadrées à
la formule naïve~\eqref{eq:eq5}, évaluée littéralement sur \emph{toutes} les cases d'un
problème jouet de taille $20\times15$ - assez petit pour que le calcul brut soit
possible. L'écart maximal est de l'ordre de $10^{-16}$, c'est-à-dire la précision d'un
nombre flottant. La mise en cache est donc une réécriture exacte, pas une
approximation.

\paragraph{Suivre la convergence sans tout recalculer.} Le même problème se pose pour
vérifier que le modèle progresse : évaluer~\eqref{eq:obj7} directement coûterait à
nouveau $O(MNK)$. La même astuce s'applique (Éq.~14 de l'article) :
\begin{equation}\label{eq:objfast}
\sum_{u}\sum_{i\notin\mathcal{R}_u} c_i\hat r_{ui}^2
=\sum_{u=1}^{M}\mathbf p_u^{\top}S^q\mathbf p_u-\sum_{(u,i)\in\mathcal{R}}c_i\hat r_{ui}^2,
\end{equation}
ce qui ramène le coût à $O(|\mathcal{R}|+MK^2)$, donc au même ordre qu'une itération.
C'est ce que calcule \texttt{eals\_local.objective}, et c'est ainsi que nous vérifions
la convergence.

\subsection{Complexité}

\begin{table}[h]
\centering\small
\caption{Coût d'une itération. $|\mathcal{R}|$ est le nombre d'interactions observées,
$M$ le nombre d'utilisateurs, $N$ celui d'articles et $K$ celui de
facteurs.}\label{tab:complexity}
\begin{tabular}{@{}lll@{}}
\toprule
Méthode & Temps par itération & Poids des cases vides \\
\midrule
ALS, Hu et al.~\cite{hu2008implicit} & $O((M+N)K^3+|\mathcal{R}|K^2)$ & uniforme seulement \\
Formule naïve~\eqref{eq:eq5}         & $O(MNK)$                       & quelconque \\
\textbf{eALS}~\cite{he2016eals}      & $O((M+N)K^2+|\mathcal{R}|K)$   & un poids $c_i$ par article \\
\bottomrule
\end{tabular}
\end{table}

Le décompte est direct. Pour un utilisateur donné et un facteur donné, le terme faisant
intervenir le cache coûte $O(K)$ et les sommes sur les interactions observées coûtent
$O(|\mathcal{R}_u|)$. En sommant sur les $K$ facteurs et les $M$ utilisateurs, une étape
revient à $O(MK^2+|\mathcal{R}|K)$, et le calcul des caches ajoute $O((M+N)K^2)$.

La conclusion pratique tient en une phrase : eALS est un facteur $K$ moins cher qu'ALS,
et son coût doit croître comme le \emph{carré} du nombre de facteurs plutôt que comme
son \emph{cube}. C'est une prédiction vérifiable, et la Section~\ref{sec:factors} la
met à l'épreuve.

\subsection{Distribuer une itération}\label{sec:algodist}

\begin{figure}[h]
\centering
\begin{minipage}{0.97\textwidth}
\begin{lstlisting}[language=,basicstyle=\ttfamily\small,numbers=none,frame=single,
                   escapechar=@,caption={},captionpos=b,label={lst:alg}]
entrée : P, Q sur le driver ; c diffuse une fois pour toutes
         U_1..U_B partitionnent les utilisateurs, I_1..I_B partitionnent les items

 0  avant la boucle :  S^q <- somme_i c_i q_i q_i^T  ;  bcP <- broadcast(P)
    (les seules fois ou ces deux objets sont calcules hors des taches)

 1  bcQ <- broadcast(Q)
 2  pour chaque bloc b, EN PARALLELE :                      # etape utilisateur
 3      P_b <- lignes U_b de bcP
 4      pour f = 1..K :  mettre a jour p_uf pour TOUS les u de U_b  -- Eq.(12), vectorise sur u
 5      emettre (U_b, P_b, P_b^T P_b)
 6  collecter ; disperser les P_b dans P ;  S^p <- somme_b P_b^T P_b     # cache, gratuit
 7  bcP <- broadcast(P)
 8  pour chaque bloc b, EN PARALLELE :                      # etape item
 9      Q_b <- lignes I_b de bcQ
10      pour f = 1..K :  mettre a jour q_if pour TOUS les i de I_b  -- Eq.(13), vectorise sur i
11      emettre (I_b, Q_b, somme_{i in I_b} c_i q_i q_i^T)
12  collecter ; disperser les Q_b dans Q ;  S^q <- somme_b (.)  # cache pour l'iteration suivante
\end{lstlisting}
\end{minipage}
\caption{Une itération eALS distribuée (\texttt{eals\_rdd.fit}). Les lignes 5 et 11
sont ce qui fait que les deux caches $S$ de l'Algorithme~1 de l'article ne coûtent
rien : chaque tâche de mise à jour renvoie la matrice de Gram des lignes qu'elle vient
d'écrire.}\label{alg:dist}
\end{figure}

Trois choix de conception méritent d'être détaillés.

\paragraph{La parallélisation ne perd aucune précision.} C'est l'affirmation centrale
de l'article, et elle se justifie simplement. Pendant l'étape utilisateur, $Q$ et le
cache $S^q$ ne sont que \emph{lus} ; les seules valeurs \emph{écrites} par un bloc sont
les facteurs de ses propres utilisateurs, auxquels aucun autre bloc ne touche. Il n'y a
donc jamais deux machines qui modifient la même case.

Le résultat distribué est par conséquent identique au résultat séquentiel, à l'ordre
d'addition des nombres flottants près. C'est une propriété rare : la descente de
gradient stochastique distribuée, par exemple, ne l'a pas, car des mises à jour
concurrentes s'y écrasent mutuellement. La Section~\ref{sec:correctness} vérifie
l'affirmation numériquement.

\paragraph{Les caches sont obtenus gratuitement.} Une implémentation naïve recalculerait
$S^q$ et $S^p$ en parcourant à nouveau toutes les lignes de facteurs, soit deux passes
supplémentaires par itération. Nous procédons autrement : chaque tâche renvoie, en même
temps que ses facteurs mis à jour, la contribution partielle des lignes qu'elle vient
d'écrire (lignes 5 et 11 du pseudo-code). Le driver n'a plus qu'à additionner ces
morceaux. Le cache est donc un sous-produit du calcul, et le surcoût réseau est
négligeable : quelques matrices $K\times K$, soit environ 131\,ko par étape.

\paragraph{Le parallélisme est entre utilisateurs, pas entre facteurs.} La boucle sur
les $K$ facteurs ne peut pas être parallélisée : le facteur $f+1$ doit voir la valeur
mise à jour du facteur $f$, sinon on ne résout plus le bon problème. En revanche, rien
n'empêche de traiter le facteur $f$ de \emph{tous} les utilisateurs d'un bloc en même
temps, puisqu'ils sont indépendants. C'est exactement ce que fait notre noyau, et cela
change l'échelle du problème : la boucle Python s'exécute $K$ fois par étape au lieu de
$M\times K$ fois. La Section~\ref{sec:code} montre le code correspondant.

%==========================================================================
\section{Commentaire sur les principaux fragments de code}\label{sec:code}
%==========================================================================

Cette section commente les quelques fragments de code où se joue quelque chose de non
trivial. Le code complet figure en Annexe~\ref{app:code}.

\subsection{Le noyau de calcul}

C'est le cœur du projet : une vingtaine de lignes qui appliquent la règle~\eqref{eq:pu}
non pas à un utilisateur, mais à un bloc entier d'utilisateurs d'un seul coup.

\begin{lstlisting}[caption={\texttt{eals\_local.update\_users} - Éq.~\eqref{eq:pu},
vectorisée sur les utilisateurs d'un bloc.},label={lst:kernel}]
def update_users(PT, QT, idx, rowid, cvec, Sq, lam, w=1.0, rhat=None):
    K, m = PT.shape
    c_e = cvec[idx]
    wmc = w - c_e                       # (w_ui - c_i)
    wr  = w * 1.0                       # w_ui * r_ui, r_ui = 1
    if rhat is None:
        rhat = _predict(PT, QT, rowid, idx)
    for f in range(K):
        qf = QT[f][idx]
        rhat -= PT[f][rowid] * qf       # r_hat^f = r_hat - p_uf q_if
        num = np.bincount(rowid, (wr - wmc * rhat) * qf, minlength=m)
        den = np.bincount(rowid, wmc * qf * qf,          minlength=m)
        num -= Sq[:, f] @ PT - PT[f] * Sq[f, f]     # - sum_{k!=f} p_uk s^q_kf
        den += Sq[f, f] + lam
        PT[f] = num / den
        rhat += PT[f][rowid] * qf
    return rhat
\end{lstlisting}

\begin{itemize}\itemsep3pt
\item \textbf{Le code suit la formule terme à terme.} \texttt{num} est le numérateur
de~\eqref{eq:pu} et \texttt{den} son dénominateur. La ligne
\texttt{Sq[:,f] @ PT - PT[f]*Sq[f,f]} calcule la somme sur les facteurs autres que $f$ :
on somme sur tous les facteurs, puis on retranche le terme en trop, ce qui est plus
rapide que de construire une matrice trouée.
\item \textbf{\texttt{np.bincount} fait tout le travail.} C'est la fonction qui permet
de calculer, en une seule passe, une somme séparée pour chaque utilisateur du bloc.
C'est elle qui rend possible le traitement simultané de tout un bloc, au lieu d'une
boucle Python par utilisateur.
\item \textbf{Les facteurs sont stockés transposés.} La boucle interne accède à un
facteur entier à la fois. En rangeant la matrice dans l'autre sens, chaque accès serait
dispersé en mémoire au lieu d'être contigu, ce qui coûte cher sur des tableaux de cette
taille.
\item \textbf{Les prédictions sont mises à jour, jamais recalculées.} Le tableau
\texttt{rhat} contient les prédictions courantes. À chaque facteur, on retire sa
contribution (\texttt{rhat -= ...}), on l'utilise, puis on rajoute la nouvelle
(\texttt{rhat += ...}). Sans cette mise à jour incrémentale, il faudrait recalculer
chaque prédiction depuis zéro, ce qui multiplierait le coût par $K$.
\item \textbf{Aucun grand tableau intermédiaire n'est créé.} Rassembler d'un coup les
facteurs de toutes les interactions demanderait 2,5\,Go sur ml-10m. Nous récupérons une
colonne à la fois, ce qui ramène le besoin mémoire au nombre d'interactions, sans
changer le volume total de calcul.
\end{itemize}

\subsection{Construire les blocs une seule fois}

\begin{lstlisting}[caption={\texttt{eals\_rdd.build\_blocks} - un élément RDD par
partition, contenant le bloc sous forme de tableaux NumPy.},label={lst:blocks}]
def build_blocks(df, key, val, n_parts):
    def to_arrays(it):
        ks, vs = [], []
        for k, v in it:
            ks.append(k); vs.append(v)
        ks = np.asarray(ks, dtype=np.int64); vs = np.asarray(vs, dtype=np.int64)
        ids, rowid = np.unique(ks, return_inverse=True)
        o = np.argsort(rowid, kind="stable")
        return iter([(ids, rowid[o].astype(np.int64), vs[o])])

    rdd = df.select(key, val).rdd.map(lambda r: (r[0], r[1]))
    return rdd.partitionBy(n_parts).mapPartitions(to_arrays, preservesPartitioning=True)
\end{lstlisting}

\begin{itemize}\itemsep3pt
\item \textbf{Les données ne sont réparties qu'une seule fois.} On regroupe les
interactions par utilisateur une fois pour toutes, et on garde le résultat en mémoire.
Refaire ce regroupement à chaque itération obligerait Spark à redéplacer toutes les
interactions sur le réseau à chaque fois. Sur ml-10m, c'est la différence entre une
itération de 4 secondes et une de 40.
\item \textbf{Un bloc entier tient dans un seul objet.} Chaque partition est réduite à
trois tableaux NumPy, plutôt qu'à un objet Python par utilisateur. Sur ml-10m, la
seconde option ferait payer le coût des objets Python 70\,000 fois par étape, contre une
poignée de fois ici.
\item \textbf{Les interactions sont triées à la construction.} Le tri rend les accès
mémoire séquentiels dans la boucle de calcul, au lieu de sauter au hasard dans le
tableau des facteurs.
\item \textbf{Les utilisateurs les plus actifs sont répartis entre les blocs.} Les
identifiants sont attribués par fréquence décroissante, de sorte que la répartition par
défaut disperse les gros utilisateurs au lieu de les entasser dans le même bloc.
L'enjeu est réel : sur ml-10m, l'utilisateur le plus actif a 7\,063 interactions contre
68 pour l'utilisateur médian. Un bloc qui en concentrerait plusieurs ferait attendre
tous les autres.
\end{itemize}

\subsection{Diffusion du modèle et localisation de l'état}

\begin{lstlisting}[caption={\texttt{eals\_rdd.fit} - le corps d'une itération.},label={lst:iter}]
bcQ = sc.broadcast(QT)
res = ublocks.map(lambda b: _user_task(b, bcP, bcQ, bcC, lam, w, Squ)).collect()
Sp = np.zeros((K, K))
for ids, blk, sp in res:
    PT[:, ids] = blk          # disperser le bloc dans la matrice de facteurs complète
    Sp += sp                  # ... et accumuler le cache S^p gratuitement
bcP.unpersist(); bcP = sc.broadcast(PT)
res = iblocks.map(lambda b: _item_task(b, bcP, bcQ, bcC, lam, w, Sp)).collect()
\end{lstlisting}

\begin{itemize}\itemsep3pt
\item \textbf{Deux diffusions par itération, pas quatre.} La matrice $Q$ est envoyée une
fois en début d'itération et sert aux \emph{deux} étapes ; $P$ est envoyée après l'étape
utilisateur et resservira à l'itération suivante. Le détail compte, car chaque machine
doit décompresser sa copie : une diffusion inutile se paie autant de fois qu'il y a de
processus.
\item \textbf{Libérer explicitement les copies diffusées.} Spark ne récupère pas
automatiquement la mémoire d'une variable diffusée lorsqu'on cesse de s'en servir. Sans
l'appel explicite à \texttt{unpersist}, le driver accumule une copie de $P$ par
itération et finit par saturer.
\item \textbf{Le modèle ne vit pas dans Spark.} Les facteurs sont un simple tableau
NumPy côté driver ; Spark ne sert qu'à distribuer le calcul, pas à stocker l'état. C'est
un choix délibéré, et il évite le piège décrit au fragment suivant.
\item \textbf{C'est aussi la limite de la conception.} Le modèle entier doit tenir sur
le driver, et chaque processus de calcul doit pouvoir en garder une copie. À $10^8$
utilisateurs et $K=128$, cela ferait 100\,Go et l'approche s'effondrerait. La réponse à
cette échelle serait le schéma par jointures évoqué en Section~\ref{sec:solution}.
\end{itemize}

\subsection{Croissance du plan d'exécution dans une boucle Spark}\label{sec:lineage}

Spark ne calcule rien tant qu'on ne lui demande pas de résultat : il se contente
d'enregistrer la recette des opérations. C'est efficace, sauf dans une boucle, où la
recette s'allonge à chaque tour. La boucle d'entraînement y échappe, puisque le modèle
n'y est pas stocké. La préparation des données, elle, n'y échappe pas : le filtrage
10-core est une boucle de jointures, et c'est là que le problème s'est manifesté.

\begin{lstlisting}[caption={\texttt{data.kcore} - la ligne qui permet à une boucle de
jointures Spark de survivre.},label={lst:kcore}]
nxt = df.join(ok_u, "user").join(ok_i, "item").select("user", "item", "ts")
# localCheckpoint, pas cache : cache() garde le plan logique vivant, donc apres ~8
# tours le plan est assez profond pour que Spark meure en *affichant* le plan
# (TreeNode.generateTreeString). Le checkpoint tronque la lignee en un simple scan
# des lignes materialisees.
nxt = nxt.localCheckpoint(eager=True)
df.unpersist()
df = nxt
\end{lstlisting}

\begin{itemize}\itemsep3pt
\item \textbf{Mettre en cache ne suffit pas.} On pourrait croire qu'un \texttt{cache()}
règle le problème. Il matérialise bien les \emph{données}, mais Spark conserve la
recette complète, au cas où il devrait tout recalculer. Elle continue donc de
s'allonger, et le driver finit par passer plus de temps à préparer le calcul que les
machines à l'exécuter. \texttt{localCheckpoint} coupe court en remplaçant la recette par
une simple lecture des données déjà calculées.
\item \textbf{Le symptôme est trompeur.} Sur Amazon Movies, après environ huit tours de
filtrage, le driver s'arrêtait sur une erreur de mémoire saturée. Or ce n'est pas en
traitant les données que Spark manquait de mémoire, mais en \emph{affichant} la recette
devenue gigantesque. Cela ressemble à un problème de volume de données ; c'en est un de
taille de plan d'exécution.
\item \textbf{La leçon générale.} Dans une boucle Spark, il faut périodiquement couper
l'historique des opérations. Notre boucle d'entraînement y échappe seulement parce que
le modèle vit en dehors de Spark, ce qui est précisément le choix de la
Section~\ref{sec:algodist}.
\end{itemize}

\subsection{Environnement d'exécution des processus distants}

Le notebook de l'Annexe~\ref{app:nb} importe les modules en ajoutant \texttt{../src}
à \texttt{sys.path}. Cela fonctionne sur le driver et échoue sur les workers :

\begin{lstlisting}[caption={\texttt{spark\_utils.get\_spark} - trois lignes sans
lesquelles rien ne fonctionne en dehors de \texttt{src/}.}]
src = os.path.dirname(os.path.abspath(__file__))
if src not in os.environ.get("PYTHONPATH", "").split(os.pathsep):
    os.environ["PYTHONPATH"] = os.pathsep.join(
        [src] + [p for p in os.environ.get("PYTHONPATH", "").split(os.pathsep) if p])
\end{lstlisting}

Les processus de calcul lancés par Spark sont des processus Python distincts. Ils
n'héritent pas des chemins d'import du programme principal : tout code envoyé sur un
worker et qui importe nos modules échoue donc avec un \texttt{ModuleNotFoundError}.
L'erreur est d'autant plus déroutante qu'elle survient à l'intérieur d'une tâche
distante et remonte enfouie sous plusieurs traces de pile Java.

La correction tient en trois lignes : les processus héritent en revanche des variables
d'environnement, il suffit donc d'y inscrire le chemin du code avant de démarrer Spark.
Nos scripts ne rencontraient jamais le problème, car ils s'exécutent depuis le dossier
\texttt{src/} ; le notebook, lancé depuis un autre dossier, l'a rencontré
immédiatement. C'est typiquement le genre de bug qu'un environnement de développement
unique dissimule jusqu'au jour où quelqu'un d'autre exécute le code.

\subsection{Évaluer sans construire la matrice des scores}

Pour évaluer le modèle, il faut classer l'article de test de chaque utilisateur parmi
\emph{tout} le catalogue. Calculer tous ces scores d'un coup demanderait 5,8\,Go sur
Amazon Movies. Les calculer un utilisateur à la fois en Python prendrait des minutes.
Nous procédons par blocs de 256 utilisateurs, ce qui ramène le tout à une seule
multiplication de matrices par bloc :

\begin{lstlisting}[caption={\texttt{evaluate.\_block\_eval} - classement sur le
catalogue entier, par blocs.}]
S = P[us] @ Q.T                       # (256, N) - un seul GEMM, pas 256 produits scalaires
for b, r in enumerate(buf):
    S[b, np.asarray(r[1])] = -np.inf  # masquer les items d'entrainement de cet utilisateur
st   = S[np.arange(len(buf)), tis].copy()
rank = (S > st[:, None]).sum(1) + 1   # rang (base 1) de l'item verite-terrain
\end{lstlisting}

La taille du bloc plafonne la mémoire à 45\,Mo au lieu de 5,8\,Go, et l'évaluation est
elle aussi distribuée : les utilisateurs sont répartis entre les cœurs comme le reste du
calcul.


%==========================================================================
\section{Analyse expérimentale}\label{sec:exp}
%==========================================================================

Chaque chiffre ci-dessous est produit par \texttt{src/experiments.py}, stocké en CSV
brut dans \texttt{results/}, puis transformé en figure par \texttt{src/plots.py} et en
tableau LaTeX par \texttt{report/tables.py}. Aucun chiffre de ce rapport n'est saisi à
la main. Le matériel et les deux pièges de configuration sont décrits en
Section~\ref{sec:setup} ; sauf mention contraire $K=32$, $\lambda=0,01$, $c_0=512$,
$\alpha=0,4$, graine 42, et 3 répétitions de 4 itérations dont la première est
écartée comme échauffement. La ligne de commande exacte de chaque exécution figure
dans \texttt{scripts/run\_all.sh}.

\paragraph{Avertissement sur la comparabilité des mesures.} Toute la campagne est une
séquence de jobs à 8 cœurs sur un ordinateur portable, et un portable
sous charge soutenue ne maintient pas une horloge constante. Nous avons mesuré la
conséquence plutôt que de la balayer d'un revers de main : la même configuration
ré-exécutée une heure plus tard, après un voisin plus lourd, est ressortie jusqu'à
deux fois plus lente. Nous attribuons cela à une combinaison de limitation thermique
(\emph{throttling}) et de workers JVM/Python froids dans chaque \texttt{SparkSession}
nouvellement créée, et nous n'avons pas cherché à séparer les deux effets. La règle
pratique que nous avons suivie, et que le lecteur devrait appliquer en lisant les
tableaux : \textbf{chaque conclusion ci-dessous est tirée d'un rapport calculé
\emph{à l'intérieur} d'un seul appel d'expérience} - une accélération, un facteur de
croissance, un rapport entre méthodes - et jamais de secondes absolues comparées
entre deux appels. Là où une expérience aurait sinon été exposée à cette dérive, elle
a été ré-exécutée avec les répétitions entrelacées entre configurations et un temps de
repos entre elles (\texttt{--interleave}, \texttt{--cooldown}).

%--------------------------
\subsection{Exactitude de la parallélisation}\label{sec:correctness}
%--------------------------

Rien d'autre dans cette section n'a de sens si le code distribué ne calcule pas eALS.
Deux vérifications sont effectuées, dans cet ordre. D'abord, les règles de mise à jour
elles-mêmes sont comparées à une évaluation littérale en $O(MNK)$ de
l'Éq.~\eqref{eq:eq5} sur une matrice jouet, qui concorde à $3\cdot10^{-16}$ près
(Section~\ref{sec:algo}). Ensuite - le tableau ci-dessous - l'implémentation
distribuée est comparée à la référence séquentielle. La Section~\ref{sec:algo} a
argumenté que le découpage en blocs ne peut pas changer le résultat ; ici nous le
vérifions. L'implémentation RDD est exécutée sur ml-100k avec la même graine, $K=8$ et
5 itérations, pour 1, 2, 4 et 8 blocs, et comparée entrée par entrée à la référence
NumPy séquentielle.

\begin{table}[htbp]
\centering\small
\caption{Contrôle d'exactitude (\texttt{check\_correctness.py}, ml-100k, $K=8$, 5
itérations). « blocs » est le nombre de partitions Spark. L'objectif est
l'Éq.~\eqref{eq:obj7} évaluée via l'Éq.~\eqref{eq:objfast}.}\label{tab:correctness}
\begin{tabular}{@{}lrrrrr@{}}
\toprule
Implémentation & blocs & $\max|\Delta P|$ & $\max|\Delta Q|$ & Objectif $J$ & err.\ rel. \\
\midrule
\inputrows{tables/tab_correctness.tex}
\bottomrule
\end{tabular}
\end{table}

Deux observations. Avec un seul bloc, l'exécution Spark est \textbf{identique au bit
près} à la référence NumPy - la même arithmétique dans le même ordre. Avec 2, 4 ou 8
blocs, l'écart est de $\CorrWorst$, soit quelques unités du dernier chiffre d'un
double : la seule chose qui change est l'ordre dans lequel les sommes partielles de
$S^p$ et $S^q$ sont accumulées. L'objectif concorde à une erreur relative inférieure à
$10^{-15}$. C'est la forme empirique de l'affirmation de l'article selon laquelle eALS
est trivialement parallélisable \emph{sans perte d'approximation}, et cela signifie
aussi que chaque mesure de temps ci-dessous mesure le même calcul.

%--------------------------
\subsection{Scalabilité}\label{sec:scal}
%--------------------------

\subsubsection{Scalabilité forte}\label{sec:strong}

Données fixes, ressources croissantes : \texttt{local[1]}, \texttt{local[2]},
\texttt{local[4]}, \texttt{local[8]}, avec un nombre de partitions maintenu
proportionnel au nombre de cœurs pour que le travail par tâche reste constant.

\begin{figure}[htbp]\centering
\figplot{strong_scaling}{0.98\textwidth}
\caption{Scalabilité forte sur ml-1m ($10^6$ interactions) et ml-10m ($10^7$). Les
barres d'erreur sont le min--max sur 9 échantillons par itération. Pointillés :
l'accélération linéaire idéale ; tirets : la courbe d'Amdahl ajustée à chaque jeu de
données.}\label{fig:strong}
\end{figure}

\begin{table}[htbp]
\centering\small
\caption{Scalabilité forte, $K=32$, partitions $=2\times$ cœurs.}\label{tab:strong}
\begin{tabular}{@{}lrrrr@{}}
\toprule
Cœurs $p$ & médiane s/itér. & min--max & accélération $S(p)$ & efficacité $E(p)$ \\
\midrule
\inputrows{tables/tab_strong.tex}
\bottomrule
\end{tabular}
\end{table}

\paragraph{Méthodologie.} Le nombre de blocs croît avec le nombre de cœurs ($B=2p$),
de sorte que le travail par tâche reste constant - ce que l'on ferait en production.
La contrepartie, qui compte pour la lecture des chiffres, est que la surcharge totale
par tâche croît elle aussi avec $p$ : à $p=8$ l'exécution paie seize surcharges de
bloc contre deux à $p=1$, ce qui \emph{sous-estime} l'accélération. Les répétitions
sont \emph{entrelacées} - la boucle externe est la répétition, la boucle interne le
nombre de cœurs - afin que la lente dérive thermique d'un portable sous une campagne
prolongée ne retombe pas entièrement sur la configuration exécutée en dernier. Nous
avons constaté que cela comptait : dans une première exécution non entrelacée, le point
à 8 cœurs a produit un unique échantillon de 19,7\,s contre une médiane de 8,8\,s.

\paragraph{Interprétation de l'ajustement d'Amdahl.} En suivant Amdahl~\cite{amdahl1967},
nous ajustons $T(p)=a+b/p$ par moindres carrés et rapportons $s=a/(a+b)$, la fraction
d'une itération à un cœur qui ne bénéficie pas de cœurs supplémentaires. Ajuster le
modèle directement est plus robuste que de lire $s$ sur le point à 8 cœurs, qui est
exactement le point le plus exposé à la contention mémoire.

\paragraph{Résultats.} L'exécution de référence - ml-10m, répétitions
entrelacées - atteint $S(8)=\SpeedupBig\times$ sur un possible $8\times$, une
efficacité de $26\,\%$, avec une fraction série d'Amdahl ajustée de $\SerialBig$. Sur
ml-1m, où les données sont dix fois plus petites et où le coût fixe d'une itération en
représente une part bien plus grande, $S(8)=\SpeedupSmall\times$ et $s=\SerialSmall$.
Les deux courbes s'aplatissent entre 4 et 8 cœurs : le quatrième cœur rapporte encore
quelque chose, le huitième pratiquement pas.

\paragraph{Origine de la fraction non parallélisable.} Quatre choses ne se réduisent pas quand $p$
croît :
\begin{itemize}\itemsep2pt
\item \textbf{La désérialisation des broadcasts}, payée une fois par worker Python et
par broadcast, et donc \emph{croissante} avec $p$. C'est le terme dominant, et il est
invisible dans toute décomposition mesurée côté driver : \texttt{sc.broadcast} n'écrit
le tableau qu'une fois, tandis que la moitié coûteuse - la désérialisation - se
produit dans chaque worker Python et est imputée à la tâche, pas au driver.
\item \textbf{La colle côté driver} : créer les broadcasts, disperser $2\times$ $B$
blocs dans $P$ et $Q$, réduire les partiels de Gram.
\item \textbf{La soumission des jobs et l'ordonnancement des tâches}, $2B$ tâches par
itération.
\item \textbf{La bande passante mémoire}, qui n'est pas du tout du travail série mais
se comporte comme tel : le noyau est un flux de gathers et de sommes par segment sur
des tableaux bien plus grands que le L2, et les huit cœurs d'un portable série M
partagent un seul contrôleur mémoire. C'est pourquoi le plafond est une propriété de
la machine plutôt que du code, et pourquoi un cluster de huit nœuds à un cœur
paraîtrait meilleur qu'un seul nœud à huit cœurs.
\end{itemize}

\subsubsection{Scalabilité en nombre de facteurs latents}\label{sec:factors}

C'est l'expérience qui décide si l'affirmation centrale de l'article survit au contact
d'une implémentation réelle. Par itération, eALS coûte
$O\big((M{+}N)K^2+|\mathcal{R}|K\big)$ et ALS $O\big((M{+}N)K^3+|\mathcal{R}|K^2\big)$.
Lire cela comme « $K^2$ contre $K^3$ » est une erreur, et l'erreur se manifeste
immédiatement dans les mesures. Les deux bornes sont la \emph{somme} de deux termes,
et lequel domine dépend de la forme des données :
\begin{equation}\label{eq:kstar}
(M{+}N)K^2 \;\ge\; 2|\mathcal{R}|K
\quad\Longleftrightarrow\quad
K \;\ge\; K^{\star}=\frac{2|\mathcal{R}|}{M+N}.
\end{equation}
Sous $K^{\star}$, eALS est essentiellement \emph{linéaire} en $K$ ; ce n'est
qu'au-dessus que le terme quadratique prend le dessus. $K^{\star}=\CrossKMl$ sur ml-1m
et $K^{\star}=\CrossKAmz$ sur Amazon Movies, donc les deux jeux de données se situent
de part et d'autre de la frontière intéressante et nous exécutons le balayage sur les
deux. (Pour ALS, le même raisonnement donne un régime cubique au-dessus de
$|\mathcal{R}|/(M{+}N)$, c.-à-d.\ la moitié de $K^{\star}$.)

\begin{figure}[htbp]\centering
\figplot{factor_scaling}{0.92\textwidth}
\caption{Coût d'une itération en fonction de $K$, 8 cœurs, log-log. Les repères gris
sont $K$, $K^2$ et $K^3$ purs, ancrés au plus petit $K$. « plage haute » est la pente
ajustée sur les deux dernières octaves seulement, où le coût constant de Spark ne
compte plus.}\label{fig:factors}
\end{figure}

\begin{table}[htbp]
\centering\small
\caption{Secondes par itération en fonction de $K$, 8 cœurs. eALS est notre
implémentation Spark RDD ; ALS est
\texttt{pyspark.ml.recommendation.ALS(implicitPrefs=True)}, c.-à-d.\ Hu et
al.~\cite{hu2008implicit}, chronométré par son coût \emph{marginal}
$(T_{11}-T_1)/10$ afin d'exclure le coût fixe de construction de ses blocs. ALS n'a pas
été exécuté au-delà de $K=128$ : une seule itération y coûte déjà des
minutes.}\label{tab:factors}
\begin{tabular}{@{}lrrr@{}}
\toprule
$K$ & eALS (s/itér.) & MLlib ALS (s/itér.) & rapport \\
\midrule
\inputrows{tables/tab_factors.tex}
\bottomrule
\end{tabular}
\end{table}

\paragraph{Résultats.} L'affirmation la plus nette est un facteur de
croissance sur une fenêtre fixe plutôt qu'un exposant ajusté. Sur \textbf{Amazon
Movies}, entre $K=32$ et $K=128$, le coût d'une itération eALS est multiplié par
$\GrowthEalsAmz$ et celui d'une itération MLlib ALS par $\GrowthAlsAmz$ - ALS croît
près de cinq fois plus vite sur la même fenêtre, ce qui est l'affirmation qualitative
de l'article. La conséquence en termes absolus est un point de croisement : à
$K=\KAtRatioMinAmz$, MLlib ALS est $\RatioAtMinAmz\times$ \emph{plus rapide} que notre
eALS, et à $K=\KAtRatioAmz$, eALS est $\RatioAtMaxAmz\times$ plus rapide. Sur
\textbf{ml-1m}, la même fenêtre donne $\GrowthEalsMl$ contre $\GrowthAlsMl$ : le
rapport entre les deux taux de croissance est essentiellement le même, sur un jeu de
données dont la forme est complètement différente.

\paragraph{La forme théorique s'ajuste, et elle nous dit où se trouve réellement le
croisement.} Ajuster $t = a + bK + cK^2$ pour eALS et $t = a + bK^2 + cK^3$ pour ALS
- les deux bornes de complexité, avec des constantes libres - reproduit les
mesures presque exactement : $R^2 = \FitRtwoEalsAmz$ et $\FitRtwoAlsAmz$ sur Amazon
Movies, $\FitRtwoEalsMl$ et $\FitRtwoAlsMl$ sur ml-1m. Les constantes additives sont
$\FitConstEalsAmz$\,s pour eALS contre $\FitConstAlsAmz$\,s pour MLlib : notre
implémentation Python paie environ le double du coût fixe par itération du code JVM
natif, ce qui explique qu'ALS gagne à petit $K$.

Le chiffre intéressant est le rapport $b/c$, le $K$ auquel le terme d'ordre supérieur
dépasse le terme inférieur. Le décompte de flops prédit
$K^{\star}=2|\mathcal{R}|/(M{+}N)=\CrossKAmz$ sur Amazon Movies ; le croisement
\emph{ajusté} est $\FitCrossEalsAmz$, neuf fois plus élevé. Sur ml-1m la prédiction
est $\CrossKMl$ et l'ajustement donne $\FitCrossEalsMl$.

\paragraph{Pourquoi le décompte de flops est faux d'un ordre de grandeur, et pourquoi
c'est la chose la plus utile que dit cette expérience.} Les deux termes de
$O((M{+}N)K^2+|\mathcal{R}|K)$ ont des constantes \emph{par opération} très
différentes dans toute implémentation réelle, et la borne de complexité les traite
comme interchangeables. Dans notre noyau :
\begin{itemize}\itemsep2pt
\item le terme en $K^2$ est \texttt{Sq[:, f] @ PT}, un produit matrice-vecteur dense
que BLAS exécute à plusieurs GFLOP/s hors cache ;
\item le terme en $|\mathcal{R}|K$ est \texttt{QT[f][idx]} suivi de deux sommes par
segment \texttt{np.bincount} - un gather aléatoire et un scatter aléatoire sur des
tableaux de $\textit{nnz}$ éléments, limités par la latence mémoire, un ordre de
grandeur plus lent par opération.
\end{itemize}
Un décompte de flops place donc $K^{\star}$ beaucoup trop bas, et l'exposant mesuré
reste proche de $\SlopeHiEalsAmz$ sur le haut de la plage où un décompte de flops pur
prédirait presque $2$. Ce n'est pas une erreur de l'analyse de l'article, qui est
asymptotiquement correcte ; c'est l'écart ordinaire entre un décompte d'opérations et
une hiérarchie mémoire, et cela a une conséquence pratique : \textbf{sur nos deux
jeux de données, eALS se comporte, sur toute la plage de $K$ qu'on peut se permettre
d'entraîner, bien plus près d'un algorithme $O(|\mathcal{R}|K)$ que d'un algorithme
$O((M{+}N)K^2)$} - ce qui est meilleur pour eALS que ne le suggère sa propre borne
de complexité.

\paragraph{Réserve sur la méthode de référence.} L'exposant de MLlib ($\SlopeHiAlsAmz$ sur la
plage haute) est lui aussi bien en dessous de son $3$ nominal, pour une raison du même
type plus une autre : MLlib n'utilise pas une inversion de manuel. Elle accumule une
équation normale et la factorise par Cholesky sur un stockage triangulaire compact, en
code natif. L'article fait la même observation à propos de sa propre référence ALS en
Java. La formulation honnête de cette expérience est donc « un algorithme quadratique
en $K$ écrit en Python contre un algorithme quasi-cubique bien optimisé écrit en code
natif » - et le quadratique gagne quand même à partir de $K=\FlipKAmz$, par une
marge qui croît.

\subsection{Convergence et précision}\label{sec:quality}
%--------------------------

\begin{figure}[htbp]\centering
\figplot{convergence}{0.92\textwidth}
\caption{L'objectif~\eqref{eq:obj7}, calculé via~\eqref{eq:objfast}, à chaque
itération ($K=32$, $c_0=512$, $\alpha=0,4$). La décroissance monotone est la preuve
pratique que les règles de mise à jour sont les minimiseurs de coordonnées
exacts.}\label{fig:conv}
\end{figure}

\begin{table}[htbp]
\centering\small
\caption{Convergence de l'objectif, $K=32$, 8 cœurs. « itérations à 1\,\% » est la
première itération dont l'objectif est à moins de 1\,\% de la valeur après 30
itérations.}\label{tab:convergence}
\begin{tabular}{@{}lrrrrr@{}}
\toprule
Jeu de données & $J$ après 1 itér. & $J$ après 30 & décroissance totale & itérations à 1\,\% & secondes \\
\midrule
\inputrows{tables/tab_convergence.tex}
\bottomrule
\end{tabular}
\end{table}

L'objectif décroît de façon monotone à chaque itération sur les deux jeux de données,
ce qui est le contrôle pratique que~\eqref{eq:pu} et~\eqref{eq:qi} sont réellement les
minimiseurs de coordonnées exacts - un pas de descente de coordonnées même
légèrement faux se manifesterait immédiatement ici par une remontée. La convergence
est rapide dans le sens qui compte : $\ConvItersMl$ itérations sur ml-1m et
$\ConvItersAmz$ sur Amazon Movies amènent $J$ à moins de 1\,\% de sa valeur après 30,
et l'essentiel de la décroissance totale se produit dans les trois premières. Un budget
de 20 itérations, utilisé partout ailleurs dans ce rapport, est donc largement au-delà
du point où l'objectif cesse de bouger.

\begin{table}[htbp]
\centering\small
\caption{Protocole hors-ligne : leave-one-out, classement contre le catalogue
\emph{entier}. $K=32$, $\lambda=0,01$, $c_0=512$, $\alpha=0,4$, 8 cœurs. « fit » est le
temps d'entraînement (horloge murale) du modèle rapporté.}\label{tab:quality}
\begin{tabular}{@{}lrrrrr@{}}
\toprule
Méthode & HR@10 & HR@100 & NDCG@10 & NDCG@100 & fit (s) \\
\midrule
\inputrows{tables/tab_quality.tex}
\bottomrule
\end{tabular}
\end{table}

\begin{figure}[htbp]\centering
\figplot{quality_vs_time}{0.92\textwidth}
\caption{Précision en fonction du temps d'entraînement \emph{réel} (horloge murale),
pas du numéro d'itération. Chaque marqueur est une exécution d'entraînement complète
de la durée donnée, évaluée sur l'ensemble retenu.}\label{fig:qualtime}
\end{figure}

Trois constats, sur Amazon Movies, le jeu de données de l'article.

\textbf{(i) eALS bat nettement l'ALS de MLlib en précision.} HR@100 passe de $\HRa$ à
$\HRe$, soit $+\EalsOverAls$\,\%, avec NDCG@100 s'améliorant dans la même proportion.
Les deux sont entraînés sur les mêmes données avec le même $\lambda$, le même $K$ et
le même nombre d'itérations. L'article rapporte le même ordre sur son propre
instantané Amazon, et l'attribue à l'objectif sensible à la popularité plutôt qu'à
l'optimiseur - ce que confirme le point suivant.

\textbf{(ii) La pondération sensible à la popularité vaut à elle seule
$\PopGain$\,\% de HR@100.} La seule différence entre les deux lignes eALS du
Tableau~\ref{tab:quality} est $\alpha=0,4$ contre $\alpha=0$ ; tout le reste, y
compris $c_0$, la graine et le nombre d'itérations, est identique. La variante
uniforme se situe entre MLlib ALS et eALS complet, exactement là où l'argument de
l'article dit qu'elle devrait se trouver : une partie du gain vient d'une bonne
optimisation de l'objectif sur l'ensemble des données, une partie d'une bonne
pondération des données manquantes. Les deux modèles sont très au-dessus du plancher
MostPopular ($\PopFloor\times$ son HR@100).

\textbf{(iii) Sur la métrique qui compte en pratique - la précision pour un budget
de temps donné - eALS domine à chaque budget.} La Figure~\ref{fig:qualtime} est la
comparaison honnête, car une comparaison par itération masquerait le fait qu'une
itération ALS et une itération eALS ne représentent pas la même quantité de travail.
eALS atteint la meilleure précision qu'ALS n'atteint jamais (HR@100 $=\AlsBestHR$)
après $\EalsTimeToAls$\,s d'entraînement, contre $\AlsTimeToBest$\,s pour ALS, puis
continue de s'améliorer alors qu'ALS a plafonné : MLlib ALS est plat à partir de
l'itération~8, à HR@100 $\approx 0,118$. Notez que ceci n'est \emph{pas} une
affirmation sur la vitesse brute - à $K=32$ une itération ALS sur ces données est
moins chère que la nôtre, car MLlib est du code JVM natif et nous payons la
sérialisation Python. C'est une affirmation sur l'endroit où converge chaque
algorithme.


%==========================================================================
\section{Discussion : forces et faiblesses}\label{sec:discussion}
%==========================================================================

\subsection{Points forts}

\paragraph{S1 - Quadratique, pas cubique, en $K$.} C'est l'affirmation centrale de
l'article, et la Section~\ref{sec:factors} la confirme dans la forme qui résiste à
l'examen : sur la fenêtre $K=32\rightarrow128$ sur Amazon Movies, une itération eALS
croît de $\GrowthEalsAmz\times$ et une itération ALS de $\GrowthAlsAmz\times$. Comme
la précision de la MF implicite continue de s'améliorer avec $K$, le régime où elle est
la plus performante est précisément le régime où eALS gagne par la plus grande marge
- et le croisement à $K=\FlipKAmz$ se situe en dessous du $K$ que l'on veut utiliser.

\paragraph{S2 - Pas d'inversion de matrice, donc pas de problème de
conditionnement.} Chaque mise à jour de coordonnée est une division scalaire. ALS
doit inverser $Q^{\top}W^uQ+\lambda I$ une fois par utilisateur ; à $K=256$ avec un
$Q$ mal conditionné, c'est là que vivent les ennuis numériques, et c'est aussi
pourquoi MLlib a besoin d'un chemin Cholesky avec un plancher de régularisation. eALS
n'a pas un tel objet.

Une réserve que l'article ne soulève pas, et que nous pouvons préciser. Le
dénominateur de~\eqref{eq:pu} est
$\sum_{i\in\mathcal{R}_u}(w_{ui}-c_i)q_{if}^2+s^q_{ff}+\lambda$, dans lequel
$s^q_{ff}=\sum_i c_iq_{if}^2>0$ et $\lambda>0$, mais la première somme n'est
\emph{pas} de signe défini : rien dans le modèle n'interdit $c_i>w_{ui}$. Comme
$c_i \approx c_0/N$ en moyenne, cela se produit sur les petits catalogues - sur
ml-100k avec le $c_0=512$ de l'article et $N=1\,152$, l'item le plus populaire atteint
$c_i=1,07>1=w_{ui}$. Nous n'avons jamais observé d'objectif non monotone
(Figure~\ref{fig:conv}), car $s^q_{ff}$ domine en pratique, mais la garantie est
empirique plutôt que structurelle, et c'est une raison de faire varier $c_0$ avec $N$
plutôt que de transposer une valeur d'un jeu de données à l'autre.

\paragraph{S3 - Un parallélisme exact.} La
Section~\ref{sec:correctness} mesure un écart maximal de $\CorrWorst$ entre la
référence NumPy séquentielle et les exécutions Spark à 1, 2, 4 et 8 blocs, et
l'exécution à un bloc est identique au bit près. Rien n'a besoin d'être ajusté,
aucune obsolescence (\emph{staleness}) n'a besoin d'être tolérée, aucun taux
d'apprentissage n'a besoin d'être re-recherché quand le nombre de workers change -
ce que la SGD distribuée ne peut précisément pas promettre.
Opérationnellement cela vaut plus qu'il n'y paraît : cela signifie que le nombre de
partitions est un pur bouton de performance, et qu'une expérience de scalabilité ne
peut pas changer silencieusement le modèle qu'elle mesure.

\paragraph{S4 - Pas de taux d'apprentissage.} Chaque mise à jour est le minimiseur
exact de l'objectif le long d'une coordonnée. Comparé à RCD~\cite{devooght2015rcd},
qui a besoin d'une recherche linéaire à chaque pas, ou à BPR~\cite{rendle2009bpr}, dont
la précision est visiblement fonction de la taille du pas, cela retire une dimension
entière du budget d'ajustement. Notre propre expérience le confirme : aucune
configuration que nous ayons exécutée n'a jamais divergé, et l'objectif décroît de
façon monotone à chaque itération sur les deux jeux de données
(Figure~\ref{fig:conv}).

\subsection{Points faibles}

\paragraph{W1 - Le balayage de coordonnées est intrinsèquement séquentiel à
l'intérieur d'une ligne.} La boucle en $K$ est du Gauss--Seidel : $p_{u,f+1}$ doit
voir le $p_{uf}$ mis à jour. Le parallélisme est donc limité à $\min(M,B)$ dans
l'étape utilisateur et $\min(N,B)$ dans l'étape item, et le travail par bloc contient
toujours une boucle Python de longueur $K$. C'est cette boucle qui a obligé
l'implémentation à être vectorisée \emph{entre} utilisateurs : c'est le seul axe
restant. Dans un contexte à peu d'utilisateurs et à $K$ énorme - l'opposé d'un
système de recommandation - eALS se paralléliserait mal.

\paragraph{W2 - Diffuser $Q$ ne passe pas à l'échelle des catalogues réels.} La
conception RDD exige $P$ et $Q$ sur le driver et une copie désérialisée de chacun par
worker Python. À $N=10^6$ items et $K=128$ cela fait 1\,Go par worker. C'est la limite
structurelle de la conception retenue en Section~\ref{sec:solution}, et elle est
assumée : l'alternative - faire transiter les facteurs par des jointures avec
shuffle - déplace $O(|\mathcal{R}|K)$ au lieu de $O((M+N)K)$, soit un facteur
$\approx 300$ sur ml-1m, et se paie à chaque itération. Il n'y a pas de repas gratuit
ici : à une certaine taille de catalogue, la conception fondée sur les jointures gagne
simplement parce que la conception broadcast cesse de tenir en mémoire, et la bonne
réponse industrielle n'est probablement ni l'une ni l'autre, mais le schéma à
partitionnement par blocs de MLlib, où chaque ligne de facteur n'est envoyée qu'aux
blocs qui en ont besoin.

\paragraph{W3 - Le schéma de pondération exige une recherche bidimensionnelle
coûteuse.} $c_0$ et $\alpha$ interagissent, les propres optima de l'article diffèrent
entre ses deux jeux de données ($c_0=512,\alpha=0,4$ sur Yelp ; $c_0=64,\alpha=0,5$
sur Amazon), et chaque point de la grille est un entraînement complet - une
recherche bidimensionnelle sérieuse représente donc des dizaines d'entraînements. Nous
avons repris les valeurs de l'article plutôt que de les rechercher, ce qui est
précisément la raison pour laquelle nos chiffres de précision ne doivent pas être lus
comme un optimum. Pire, la sensibilité n'est pas symétrique : un $c_0$ trop petit
détruit la précision, un $c_0$ trop grand la dégrade doucement, donc une grille
grossière peut facilement tomber sur un plateau et déclarer victoire.

\paragraph{W4 - Un protocole d'évaluation indulgent, en particulier envers la
pondération par popularité.} Trois problèmes distincts :
\begin{itemize}\itemsep2pt
\item \emph{Le leave-one-out par horodatage n'est pas un découpage réaliste.} Il
laisse fuir le futur : l'ensemble d'entraînement de l'utilisateur $u$ contient des
interactions survenues après l'interaction de test de l'utilisateur $v$. L'article le
dit lui-même et ajoute pour cette raison un protocole en ligne fondé sur une coupure
chronologique globale, que nous n'avons pas reproduit ; tous nos chiffres de précision
héritent donc de ce biais.
\item \emph{Des négatifs échantillonnés rendraient les chiffres incomparables.} Nous
classons contre le catalogue entier précisément parce qu'échantillonner 100 négatifs~\cite{krichene2020sampled}
est connu pour changer non seulement les valeurs mais l'\emph{ordre} relatif des
systèmes comparés (Krichene et Rendle, KDD 2020). Cela coûte $O(MNK)$ par évaluation et c'est
la principale raison pour laquelle notre évaluation est distribuée.
\item \emph{HR et NDCG ne mesurent rien sur la diversité ou la couverture du
catalogue.} C'est la critique que nous jugeons la plus importante, car elle pointe
vers la contribution même de l'article. Pondérer les données manquantes par la
popularité indique au modèle qu'un manque sur un blockbuster est une forte preuve de
désintérêt - mais le même mécanisme pousse les items impopulaires vers un score
prédit bas pour tout le monde, car il est bon marché de se tromper sur eux. Une
métrique qui demande seulement « l'item retenu était-il dans le top 100 » ne peut pas
voir la concentration qui en résulte, et les items retenus sont eux-mêmes distribués
selon la popularité, donc la métrique est biaisée positivement en faveur de tout ce
qui amplifie la popularité. Nous mesurons un vrai
gain à partir de $\alpha>0$ (Section~\ref{sec:quality}), et nous croyons qu'il est réel ;
nous croyons aussi qu'une évaluation rapportant la couverture dans l'espace des items
montrerait qu'une partie de ce gain est un biais de popularité plutôt qu'une
meilleure modélisation des préférences. Mesurer cela est l'extension la plus utile de
ce travail.
\end{itemize}

\paragraph{W5 - eALS optimise une perte de régression, pas une perte de
classement.} L'objectif est l'erreur quadratique sur une matrice $0/1$. C'est un bon
objectif de \emph{rappel}, ce que montrent exactement nos résultats - eALS est le
plus fort en HR@100 - et un objectif de \emph{précision} médiocre. BPR, qui
optimise directement l'ordre par paires, est le concurrent naturel aux petites
coupures ; cette asymétrie entre rappel et précision est un résultat classique de la
littérature sur la recommandation top-$N$~\cite{cremonesi2010topn}. Nous n'avons pas implémenté BPR (voir plus
bas) et ne pouvons donc pas le quantifier sur nos données ; l'article rapporte le même
schéma.

Un point connexe et plus inconfortable : Rendle et al. ont
revisité iALS - la référence même à laquelle eALS est comparé - sur des
benchmarks standards et ont trouvé qu'une fois sa régularisation et son poids sur les
données non observées correctement ajustés, il rivalise avec des méthodes bien plus
récentes. Notre propre référence ALS est celle de MLlib, avec le $\lambda=0,01$ de
l'article et l'$\alpha$ par défaut de MLlib ; nous ne l'avons \emph{pas} ajustée. Une
partie de l'écart de $\EalsOverAls$\,\% que nous mesurons est donc attribuable à un
effort d'ajustement inégal, et une lecture honnête du Tableau~\ref{tab:quality}
devrait le traiter comme une borne supérieure de l'avantage d'eALS sur un iALS bien
ajusté.

\subsection{Limites de validité des mesures}

Nous les listons car la plupart sont structurelles et ne peuvent pas être corrigées
sur un portable.

\begin{enumerate}\itemsep2pt
\item \textbf{\texttt{local[$p$]} n'est pas un cluster.} Il n'y a pas de réseau, pas
de sérialisation entre machines, pas de tâche à la traîne due à un nœud lent, pas de
tolérance aux pannes exercée. Nos chiffres \emph{surestiment} donc la performance de
la conception RDD sur un vrai cluster, où diffuser $Q$ coûte de la bande passante
réseau et non une copie mémoire, et \emph{sous-estiment} d'autant la valeur de
l'alternative co-partitionnée écartée en Section~\ref{sec:solution}.
\item \textbf{La bande passante mémoire partagée est la vraie limite, pas le nombre
de cœurs.} Le noyau est une séquence de gathers et de sommes par segment sur des
tableaux bien plus grands que le L2. Passer de 1 à 8 cœurs multiplie la demande sur un
seul contrôleur mémoire, ce qui explique une grande partie de l'écart entre
l'accélération mesurée $\SpeedupBig\times$ et l'idéal $8\times$. Un cluster de 8
machines à un cœur paraîtrait meilleur qu'une seule machine à 8 cœurs ici.
\item \textbf{La dérive thermique est mesurable, et nous l'avons attrapée.} Deux
exécutions du même balayage en $K$ sur ml-1m, à une heure d'intervalle dans une
campagne continue, différaient jusqu'à un facteur deux sur les configurations les
plus lourdes : la seconde exécution suivait une mesure ALS bien plus lourde et la
machine était chaude. Un portable sous charge soutenue à 8 cœurs ne maintient pas une
horloge constante. Nos parades sont méthodologiques plutôt que physiques -
répétitions entrelacées entre configurations plutôt que regroupées, et un temps de
repos entre configurations (\texttt{--interleave}, \texttt{--cooldown}) - et les
figures montrent des barres min--max sur 9 échantillons par itération afin que
l'étalement résiduel soit visible. Sur une machine dédiée à horloge fixe, ces
précautions seraient inutiles ; sur celle-ci, elles font la différence entre une
mesure et une anecdote.
\item \textbf{Deux familles de jeux de données, un seul domaine chacune.} Amazon
Movies et MovieLens sont tous deux de nature cinématographique, tous deux
modérément denses après filtrage 10-core. Rien ici ne dit comment eALS se comporte
sur un catalogue de $10^7$ items avec une moyenne de deux interactions chacun.
\item \textbf{Le notebook est le seul test de bout en bout.} Il n'y a pas de suite de
tests unitaires ; l'exactitude est garantie par \texttt{check\_correctness.py} (qui
compare bien à une référence par force brute) et par le fait que le notebook
s'exécute jusqu'au bout. Cela a suffi à attraper de vrais bugs - notamment un
\texttt{PYTHONPATH} manquant sur les workers - mais ce n'est pas un substitut à des
tests.
\item \textbf{Tout ce qui figurait dans l'énoncé n'a pas été mesuré.} Nous n'avons
pas implémenté BPR ni RCD : les deux sont
des références pertinentes dans l'article, mais chacune est un projet en soi et le
budget est allé à l'étude de scalabilité que l'énoncé pondère le plus. Nous le disons
plutôt que de rapporter des chiffres que nous n'avons pas produits. Nous n'avons pas
non plus utilisé Yelp, pour la raison donnée en Section~\ref{sec:dataset}.
\end{enumerate}


%==========================================================================
\section{Conclusion et travaux futurs}\label{sec:conclusion}
%==========================================================================

Nous avons implémenté eALS~\cite{he2016eals} en PySpark - une version RDD fondée sur
des broadcasts, adossée à une référence NumPy séquentielle - et vérifié que les deux
produisent le même modèle à $\CorrWorst$ près, quel que soit le nombre de blocs, ce qui
est le pendant empirique de l'affirmation de l'article selon laquelle eALS est
trivialement parallélisable \emph{sans approximation}. En plus de cela, nous avons
implémenté l'évaluation leave-one-out sur catalogue entier et une campagne
expérimentale dont la sortie brute se trouve dans \texttt{results/*.csv}.

Les résultats que nous considérons comme la substance de ce rapport :

\begin{enumerate}\itemsep2pt
\item L'affirmation de complexité tient, dans le régime pour lequel elle est
énoncée. Sur Amazon Movies, le coût d'une itération entre $K=32$ et $K=128$ croît de
$\GrowthEalsAmz\times$ pour eALS et $\GrowthAlsAmz\times$ pour l'ALS implicite de
MLlib, sur les mêmes données, la même machine et le même nombre de cœurs. eALS passe
de $\RatioAtMinAmz\times$ plus lent à $K=\KAtRatioMinAmz$ à $\RatioAtMaxAmz\times$
plus rapide à $K=\KAtRatioAmz$. Ajuster les deux formes de complexité avec des
constantes libres reproduit les mesures à $R^2>0,999$, mais place le croisement entre
le terme linéaire et le terme quadratique d'eALS à $K\approx\FitCrossEalsAmz$ plutôt
qu'au $\CrossKAmz$ prédit par le décompte de flops : le terme creux est limité par la
mémoire et le terme dense relève de BLAS, et un décompte de flops ne peut pas voir la
différence.
\item La scalabilité forte est réelle mais modeste : $\SpeedupBig\times$ sur 8 cœurs
pour ml-10m, cohérent avec une fraction série d'Amdahl de $\SerialBig$. La limite
n'est pas l'algorithme - le parallélisme est exact et le travail se découpe
proprement - mais la bande passante mémoire partagée, la désérialisation des
broadcasts dans chaque worker Python, et la surcharge Python par tâche.
\item La pondération sensible à la popularité fonctionne : à $\alpha=0,4$, HR@100
passe de $\HRu$ à $\HRe$ sur Amazon Movies, soit $\PopGain$\,\% de mieux que le
réglage uniforme $\alpha=0$, toutes choses égales par ailleurs.
\item Sur la métrique qui compte en pratique - la précision pour un budget de temps
donné - eALS atteint en $\EalsTimeToAls$\,s la meilleure précision que l'ALS de
MLlib n'atteint jamais (HR@100 $=\AlsBestHR$ après $\AlsTimeToBest$\,s), puis continue
de s'améliorer.
\end{enumerate}

\paragraph{Travaux futurs.} Quatre directions, dans l'ordre où nous les prendrions.
\emph{(i) Métriques de couverture et de biais de popularité.} L'ajout le plus utile,
de loin : mesurer la couverture du catalogue et la concentration de Gini en fonction
de $\alpha$, et voir quelle part du gain de HR relève réellement d'une meilleure
modélisation des préférences.
\emph{(ii) Une implémentation co-partitionnée.} Remplacer le broadcast de $Q$ par un
échange par blocs partitionnés où chaque ligne de facteur n'est envoyée qu'aux blocs
qui la référencent - le schéma qu'utilise l'ALS de MLlib - et ré-exécuter l'étude
de scalabilité sur un vrai cluster, où le coût du broadcast cesse d'être une copie
mémoire.
\emph{(iii) Références BPR et RCD}, pour tester l'asymétrie rappel/précision que nous
ne pouvons actuellement que citer depuis l'article.
\emph{(iv) Float32 et une boucle interne Numba ou Cython.} Le noyau est limité par la
mémoire ; diviser par deux la largeur de chaque gather devrait rapporter près d'un
facteur deux, et diviserait aussi le broadcast par deux.


%==========================================================================
\section*{Déclaration sur l'usage d'outils d'assistance}
%==========================================================================
\addcontentsline{toc}{section}{Déclaration sur l'usage d'outils d'assistance}
\noindent
La rédaction de ce rapport et l'écriture du code ont bénéficié de l'assistance d'un
modèle de langage. Les mesures rapportées proviennent de nos propres exécutions sur la
machine décrite en Section~\ref{sec:setup} ; les fichiers CSV bruts sont versionnés dans
\texttt{results/} et l'ensemble des figures, tableaux et chiffres du rapport en est
regénéré par \texttt{scripts/build\_report.sh}.

\bibliographystyle{plain}
\bibliography{refs}


\appendix
%==========================================================================
\section{Code source complet}\label{app:code}
%==========================================================================

Tout est également disponible sur \texttt{https://github.com/tomasnp/eals-spark}. Les
listings ci-dessous sont les fichiers tels quels (verbatim, en anglais dans le code
source).

\subsection{\texttt{src/spark\_utils.py} - session Spark}
\lstinputlisting{../src/spark_utils.py}

\subsection{\texttt{src/data.py} - chargement, filtrage 10-core, découpage}
\lstinputlisting{../src/data.py}

\subsection{\texttt{src/eals\_local.py} - référence NumPy et noyaux partagés}
\lstinputlisting{../src/eals_local.py}

\subsection{\texttt{src/eals\_rdd.py} - implémentation Spark RDD}
\lstinputlisting{../src/eals_rdd.py}

\subsection{\texttt{src/evaluate.py} - HR@$k$, NDCG@$k$, classement sur catalogue entier}
\lstinputlisting{../src/evaluate.py}

\subsection{\texttt{src/baselines.py} - ALS de MLlib et MostPopular}
\lstinputlisting{../src/baselines.py}

\subsection{\texttt{src/check\_correctness.py} - le contrôle d'exactitude}
\lstinputlisting{../src/check_correctness.py}

\subsection{\texttt{src/experiments.py} - la campagne de mesures}
\lstinputlisting{../src/experiments.py}

\subsection{\texttt{src/plots.py} - figures}
\lstinputlisting{../src/plots.py}

\subsection{\texttt{report/tables.py} - chaque chiffre de ce rapport}
\lstinputlisting{../report/tables.py}

\subsection{\texttt{scripts/env.sh}, \texttt{run\_all.sh}, \texttt{build\_report.sh}}
\lstinputlisting[language=bash]{../scripts/env.sh}
\lstinputlisting[language=bash]{../scripts/run_all.sh}
\lstinputlisting[language=bash]{../scripts/build_report.sh}

%==========================================================================
\section{Le notebook de démonstration}\label{app:nb}
%==========================================================================

\texttt{notebooks/demo\_small\_data.ipynb} exécute tout le pipeline de bout en bout
sur \textbf{ml-100k} (943 utilisateurs, 1\,152 items, 97\,953 interactions) en
quelques minutes sur un portable. Son but est la vérification, pas la mesure : un
lecteur devrait pouvoir exécuter « Run All » et voir, sans attendre, que chaque
affirmation de ce rapport est étayée par du code qui s'exécute. Chaque section
indique son entrée et sa sortie.

\begin{table}[h]
\centering\small
\caption{Cellules du notebook, avec entrée et sortie.}\label{tab:nb}
\begin{tabular}{@{}p{0.5cm}p{4.6cm}p{4.6cm}p{5.2cm}@{}}
\toprule
\S & Entrée & Sortie & Ce que cela démontre \\
\midrule
1 & le JDK installé & versions Java / Spark / NumPy, master Spark et parallélisme
  & l'environnement est bien celui décrit en Section~\ref{sec:setup} ; en particulier
    \texttt{JAVA\_HOME} pointe vers un JDK 17. \\
2 & \texttt{data/raw/ml-100k/u.data} (brut \texttt{user\,item\,rating\,ts}) & le
    dictionnaire de statistiques et les premières lignes du découpage d'entraînement
    & le filtrage 10-core, la ré-indexation dense et le découpage leave-one-out
    reproduisent le Tableau~\ref{tab:datasets}, ligne \texttt{ml-100k}. \\
3 & fréquences des items d'entraînement & $\min/\text{médiane}/\max$ de $c_i$ et
    $\sum_i c_i$ & l'Éq.~\eqref{eq:ci} est correctement implémentée :
    $\sum_i c_i=c_0$ pour tout $\alpha$, et $\alpha=0,4$ étale $c_i$ sur un rapport
    6:1 sur ml-100k alors que $\alpha=0$ le réduit à $c_0/N$. Cela montre aussi que
    l'item le plus populaire atteint $c_i>1$ - la réserve S2 de la
    Section~\ref{sec:discussion}. \\
4 & les paires d'entraînement sous forme de deux tableaux d'entiers & un tableau de
    $J$ par itération et sa décroissance & l'objectif~\eqref{eq:obj7} décroît de
    façon monotone : les règles de mise à jour~\eqref{eq:pu}--\eqref{eq:qi} sont
    correctes. \\
5 & le DataFrame d'entraînement, 8 partitions & temps réel par itération décomposé en
    étape utilisateur / étape item / broadcast / driver & l'anatomie d'une itération
    distribuée, et la part prise par la partie non parallèle. \\
6 & les modèles des \S4 et \S5 & $\max|\Delta P|$, $\max|\Delta Q|$ et une assertion
    (\texttt{assert}) & le contrôle d'exactitude de la
    Section~\ref{sec:correctness}, en miniature. La cellule \emph{échoue
    bruyamment} si les implémentations divergent. \\
7 & matrices de facteurs + l'ensemble de test leave-one-out & HR@10/@100 et
    NDCG@10/@100 pour MostPopular, eALS, eALS avec $\alpha=0$ et MLlib ALS & eALS bat
    le plancher de popularité avec une large marge, et $\alpha>0$ bat $\alpha=0$. \\
8 & les mêmes données réajustées avec 1, 2, 4 cœurs & secondes par itération,
    accélération, efficacité & une miniature de la Section~\ref{sec:strong}, et une
    honnête : sur ml-100k l'accélération ressort \emph{en dessous de un} - plus de
    cœurs rendent eALS plus lent, car les données sont bien trop petites pour Spark. \\
\bottomrule
\end{tabular}
\end{table}

\paragraph{Temps d'exécution.} Environ deux minutes de bout en bout sur la machine
de la Section~\ref{sec:setup}, dominées par les démarrages de session Spark
(\S8 redémarre la session une fois par nombre de cœurs, ce qui est inévitable :
\texttt{local[$p$]} est fixé à la création de la session).


\end{document}
